\chapter{Introduction}
\label{section:python-api-examples:introduction}


\begin{itemize}
  \item What is a action set
  \item How to set up a project
  \item \ldots
\end{itemize}




\subsection{Parallelisation and GPU support}

Follow the configuration recommendation of Section
\ref{section:particles:swift:hpc} before you continue.

\subsubsection{Domain decomposition}

All domain decomposition is realised via your main routine, and you don't have
to modify anything besides your main loop if you work shared memory only.
Once you enable MPI, too, you will have to make some minor modifications.
\marginpar{Not sure bout this one.}


\Peano\ allows you to control the domain decomposition via explicit splits (and
mergers) of subdomains.
However, we rely on the load balancing toolbox here.
The toolbox collects some global information and then triggers splits and merges
accordingly.
You have to perform the following steps:
\begin{itemize}
  \item Include 
\begin{code}
#include "toolbox/loadbalancing/RecursiveSubdivision.h"
\end{code}
  \item Create an instance of the load balancer on each MPI rank. It is
  convenient to insert
\begin{code}
toolbox::loadbalancing::RecursiveSubdivision  loadBalancer;
\end{code}
  just before the initial
  \texttt{tarch::mpi::Rank::getInstance().isGlobalMaster()} check. The object
  accepts a lot of configuration parameters which you can look up in the source
  code documentation.
  \item You have to inform your load balancer when an individual step has
  terminated. This tells the balancer to collect some global statistics and
  (eventually) to trigger some rebalancing.
\begin{code}
loadBalancer.finishStep();
\end{code}
  I typically use the load balancer only within the grid construction. Totally
  dynamic load balancing is something this simple pre-manufactured balancer
  doesn't properly support anyway.
\begin{code}
int gridConstructionSteps = 0;
while (peano4::parallel::SpacetreeSet::getInstance().getGridStatistics().getStationarySweeps()<5) {
  examples::particles::observers::CreateGrid createObserver;
  peano4::parallel::SpacetreeSet::getInstance().traverse(createObserver);
  gridConstructionSteps++;
  loadBalancer.finishStep();
}
\end{code}
\end{itemize}



